\documentclass[12pt,a4paper]{article}

\usepackage{fontawesome5}
\providecommand{\faPencilAlt}{\faIcon{pencil-alt}}

\usepackage{styles/couleurs}
\usepackage{styles/mise_en_page}
\usepackage{styles/boites}
\usepackage{styles/maths}
\usepackage{styles/figures}

\usepackage{enumitem}
\usepackage{array}
\usepackage{tabularx}
\usepackage{amsmath}
\usepackage{amssymb}
\usepackage{tikz}
\usepackage{xparse}

\usetikzlibrary{calc,angles,quotes,arrows.meta}

\setlist[itemize]{leftmargin=1.4em,itemsep=0.15em,topsep=0.2em}
\setlist[enumerate]{leftmargin=*,itemsep=0.22em,topsep=0.25em}

\newcommand{\lignepointillee}{\noindent\dotfill\par\vspace{0.45em}}

\newcommand{\lignesreponse}[1]{%
\begin{tcolorbox}[
enhanced,colback=white,colframe=NeutreC!40,boxrule=0.5pt,arc=1.5mm,
left=3mm,right=3mm,top=2mm,bottom=1mm,before skip=2pt,after skip=2pt]
\foreach \i in {1,...,#1}{\lignepointillee}
\end{tcolorbox}}

\newcommand{\titreseq}[1]{%
\begin{center}
{\sffamily\LARGE\bfseries\color{TitreC}Séquence 6 -- Cosinus d'un angle aigu}\\[0.2em]
{\sffamily\large\bfseries\color{SavoirC}#1}
\end{center}}

\begin{document}

\titreseq{Activité 2 -- Même angle, même coefficient}

\remarque{Les mesures sont arrondies. Les quotients obtenus ne seront donc pas
exactement identiques, mais ils doivent être très proches.}

\section*{Partie A -- Une famille d'angle \(60^\circ\)}

\renewcommand{\arraystretch}{1.55}
\begin{tabularx}{0.92\textwidth}{
|>{\centering\arraybackslash}X
|>{\centering\arraybackslash}X
|>{\centering\arraybackslash}X|}
\hline
\textbf{Hypoténuse} & \textbf{Côté adjacent} &
\(\textbf{adjacent}\div\textbf{hypoténuse}\)\\
\hline
\(4\) cm & \(2\) cm & \\
\hline
\(6\) cm & \(3\) cm & \\
\hline
\(10\) cm & \(5\) cm & \\
\hline
\end{tabularx}

Que remarques-tu ?

\lignesreponse{3}

\section*{Partie B -- D'autres familles de triangles}

Chaque colonne correspond à plusieurs triangles rectangles ayant le même angle aigu.

\renewcommand{\arraystretch}{1.55}
\begin{tabularx}{\textwidth}{
|>{\centering\arraybackslash}p{2.3cm}
|>{\centering\arraybackslash}X
|>{\centering\arraybackslash}X
|>{\centering\arraybackslash}X|}
\hline
 & \textbf{Triangle 1} & \textbf{Triangle 2} & \textbf{Triangle 3}\\
\hline
\multicolumn{4}{|c|}{\textbf{Famille d'angle \(40^\circ\)}}\\
\hline
Hypoténuse & \(5\) & \(8\) & \(12\)\\
\hline
Côté adjacent & \(3{,}83\) & \(6{,}13\) & \(9{,}19\)\\
\hline
Adjacent \(\div\) hypoténuse & & & \\
\hline
\multicolumn{4}{|c|}{\textbf{Famille d'angle \(70^\circ\)}}\\
\hline
Hypoténuse & \(5\) & \(8\) & \(12\)\\
\hline
Côté adjacent & \(1{,}71\) & \(2{,}74\) & \(4{,}10\)\\
\hline
Adjacent \(\div\) hypoténuse & & & \\
\hline
\end{tabularx}

Arrondis les quotients au centième.

\section*{Partie C -- Conjecturer}

Complète.

\begin{tcolorbox}[colback=SavoirC!5,colframe=SavoirC!60,boxrule=0.7pt,arc=2mm]
Dans des triangles rectangles qui possèdent le même angle aigu,
le quotient

\[
\text{côté adjacent}\div\text{hypoténuse}
\]

semble être \dotfill

Ce quotient dépend de \dotfill{} mais ne dépend pas de \dotfill
\end{tcolorbox}

\section*{Partie D -- Utiliser le coefficient}

Pour la famille d'angle \(40^\circ\), le coefficient est proche de \(0{,}77\).

\begin{enumerate}
  \item Quelle longueur adjacente prévoit-on avec une hypoténuse de \(20\) cm ?
  \item Quelle hypoténuse prévoit-on avec un côté adjacent de \(15{,}4\) cm ?
  \item Pourquoi le coefficient est-il inférieur à \(1\) ?
\end{enumerate}

\lignesreponse{6}

\aretenir{Le coefficient qui relie l'hypoténuse au côté adjacent
porte un nom : le \textbf{cosinus de l'angle aigu}.}

\end{document}
